\subsubsection{Volatility and Liquidity Interaction}

\begin{principle} \hlt{Volatility Is a Liquidity Price}\\
Volatility raises inventory risk, adverse-selection risk, and the probability that displayed liquidity becomes stale. Liquidity providers respond through wider spreads, smaller displayed size, faster cancellations, or withdrawal.
\end{principle}

\begin{definition} \hlt{Garman-Klass Intraday Volatility}\\
For a fixed intraday interval, the Garman-Klass estimator uses open, high, low, and close prices
\begin{equation*}
\sigma_{GK}= \sqrt{\frac{1}{2}\left(\log\frac{H}{L}\right)^2-\left(2\log 2-1\right)\left(\log\frac{C}{O}\right)^2}.
\end{equation*}
It is useful for comparing intraday uncertainty across markets when transaction-level volatility is too noisy.
\end{definition}

\begin{mechanism} \hlt{Liquidity Spiral}\\
A large order can consume depth, widen spreads, and raise volatility. Higher volatility makes liquidity providers reduce size or cancel quotes, forcing the next order to trade through less liquidity and amplifying the move.
\end{mechanism}

\begin{casestudy} \hlt{Flash Crash Mechanism}\\
The 2010 Flash Crash illustrates a liquidity spiral: a large automated futures sell program met thin liquidity, high-frequency traders recycled inventory quickly, aggressive flow migrated into cash equities, spreads widened, and some routers or venues paused.\\
The trigger was an order; the crash mechanism was fragile liquidity under linked-market stress.
\end{casestudy}

\begin{caveat} \hlt{Systemic Risk Link}\\
Fragmentation is not systemic risk by itself. It becomes dangerous when market participants react to stress in the same direction, when related venues desynchronise, or when liquidity withdrawal in one instrument propagates to hedges, ETFs, futures, and correlated stocks.
\end{caveat}

\begin{taxonomy} \hlt{Shock Propagation Channels}
\begin{enumerate}[label=\roman*.]
\setlength{\itemsep}{0pt}
\item \hlt{cross-asset hedging}: futures, ETFs, options, and cash baskets transmit inventory pressure;
\item \hlt{routing feedback}: routers chase fading liquidity or pause when fills become unreliable;
\item \hlt{quote withdrawal}: market makers widen or cancel when volatility exceeds their risk budget;
\item \hlt{reference-price dependency}: pegged and benchmark orders inherit stale or broken references;
\item \hlt{correlated constraints}: margin, risk limits, and kill switches cause simultaneous de-risking.
\end{enumerate}
\end{taxonomy}

\begin{empiricalfinding} \hlt{Trade Count as HFT Proxy}\\
Rising trades per day can proxy for higher high-frequency participation when trade sizes shrink and message intensity rises. It is not sufficient alone: the same count can reflect genuine liquidity, excessive slicing, or stress-induced churn.
\end{empiricalfinding}
